\chapter{Summary and Future Outlook}

In this dissertation, we have evaluated the properties of various complex amorphous thin film oxides synthesized by combinatorial sputtering. We mostly focused on two aspects of materials science: structure and property. Structurally, amorphous materials are inherently hard to characterize and throughout this study, we learned that this is especially true for thin film amorphous materials. However, by combining several different techniques we were able to effectively characterize their structure at both the macroscopic and local level. Taking what we learned structurally and what we learned by probing their dielectric properties, we learned something about the structure-property relations of the complex amorphous oxides.  

In chapter 3, we found that the structural and dielectric properties of two chemically similar amorphous oxides, namely Ta-Ge-O and Ta-Si-O, were significantly different. We saw that the dielectric constant of Ta-Ge-O was dramatically enhanced over linear interpolation of its end members, while Ta-Si-O did not exhibit this behavior. This type of enhancement was not observed in the refractive index for both materials. We used x-ray reflectivity to directly probe the density of the thin films as a function of composition. This technique allowed us to calculate the polarizabilities of the oxides without relying on tabulated values, avoiding any additive assumptions of the individual polarizabilities. We also explored the limitations of the Clausius-Mossotti relation applied to the two systems and found that the use of tabulated atomic polarizabilities and experimentally inferred volumes systematically underestimates the dielectric constant in both systems. This result highlights the importance of directly calculating the densities and atomic polarizabilities. We then used FTIR to explore the changes in active phonon modes of Ta-Ge-O as a function of composition. The spectra did not show any obvious features in the frequency range that would explain the anomalous dielectric behavior of Ta-Ge-O. It is likely that the phonon modes contributing to the enhancement exist at lower wavenumbers. Therefore, it might be worthwhile to conduct a careful far-infrared spectroscopy experiment to cover the sub 400 cm$^{-1}$ range in future studies. A Raman spectroscopy study may also be a good complement to FTIR. 

In chapter 4, we extended the study in chapter 3 by probing the atomic structure of pure Ta$_{2}$O$_{5}$ and Ta-Ge-O using high energy synchrotron x-ray diffraction and absorption methods. Simultaneously applying different techniques greatly reduced the uncertainty in structure space by constraining the possible structural parameters of the local environment. To reduce scattering from the substrate, we sputtered the thin films on silicon nitride windows. We found bond lengths and coordination numbers of amorphous Ta$_{2}$O$_{5}$ to be similar to values reported in previous studies, where the Ta-O coordination was found to be lower than those measured in crystalline Ta$_{2}$O$_{5}$ compounds. Past research suggests this is due to 1) Ta-O having a fairly low field strength (for example compared Si-O which is a good network former) 2) the Ta cation is large in size. 

Doping of GeO$_{2}$ in Ta$_{2}$O$_{5}$ result in substantial changes to the local environment. We identify two distinct regions in composition based on pair bond distance and coordination data. The Ta-O/Ge-O bond distance does not differ greatly from pure Ta$_{2}$O$_{5}$ for high Ta cation fractions; only when a considerable amount of GeO$_{2}$ is added does $r_{Ta-O/Ge-O}$ begin to decrease to $r_{Ge-O}$. Our EXAFS data also shows that the Ta-O structural unit experiences an anisotropic distortion due to Ge doping. We find that the most dramatic change in the local environment around Ta occurs for the metal-metal (M-M) pairs, namely Ta-Ta and Ta-Ge. Further, the trend in the O-O pair bond distance was explained as a consequence of a high Ta-Ge coordination at fairly low GeO$_{2}$ composition. When the composition of Ta and Ge cations are roughly equal, the short-range order starts to resemble that of low temperature crystalline Ta$_{2}$O$_{5}$, evidenced by the coordination number of the Ta-O inner shell. Moreover, we surmise that a major factor in the reduction of density observed in chapter 3 is due to the change in coordination of inner M-M pairs, while there is only minimal contribution from changes in M-O and M-M bond distances. This suggests that network conformation has a larger effect on the density than simple changes in bond length. \textbf{need to write a little more here}

In chapter 5, we studied the composition dependence of the dielectric properties of another complementary system, Zr-Al-O and Zr-Si-O. Specifically, we used these two systems to systematically investigate the effect of the presence or absence of crystallinity on the dielectric constant and breakdown field. An essentially linear variation of the static dielectric constant, $\varepsilon_s$, was observed as a function of composition, $x$, for compositions rich in the p-block element, i.e., $x>0.4$, for both chemical systems. However an abrupt change in $\varepsilon_s$ is found near $x \approx 0.35$, associated with the onset of crystallinity in as-deposited films. 
Measurements of the index of refraction at optical frequencies, combined with a simple Clausius-Mossotti interpretation, indicates that low-frequency (ionic) contributions to the polarizability exhibit systematic deviation with respect to values linearly interpolated from the endmembers. These trends are not consistently affected by the presence of crystalline order, but are related to changes associated with heterogeneous local oxygen coordination and bonding. There were no significant differences in the breakdown field associated with the level of crystallinity in both systems.

In chapter 6, we expand on the study in chapter 3 by substituting different cations in the Ta/Ge cation site which include the combinations Hf-Si-O, Hf-Ge-O, Ta-Sn-O, and Nb-Ge-O. Hf-Si-O and Hf-Ge-O showed similar trends in the x-ray diffraction spectra as a function of composition. Both remain amorphous until a critical Hf cation concentration is reached, upon which crystallizing into a mostly monoclinic phase mixed with possibly tetragonal and orthorhombic phases. Crystallization of the oxides were clearly reflected in the density measurements as jumps in the density for both systems. However, when we probed the dielectric constant of each system as a function composition, the two systems show dissimilar trends in $\varepsilon_{r}$, namely Hf-Si-O shows an enhanced $\varepsilon_{r}$ of approximately 28 at a concentration of 96 at\% Hf, while $\varepsilon_{r}$ for Hf-Ge-O is not enhanced for any composition. The dielectric enhancement observed for Hf-Si-O is clear indication that the monoclinic phase is not the sole phase within this composition range, since $\varepsilon_{r}$ of a pure monoclinic phase is only 16. The refractive index is linear with composition for both Hf-Si-O and Hf-Ge-O.

Ta-Sn-O and Nb-Ge-O are quite different structurally and electronically compared to Hf-Si-O and Hf-Ge-O. Inspection of the diffraction spectra indicates that both systems remain amorphous for most of the compositions probed in the study. The exception is for the Sn rich region in Ta-Sn-O ($<$ 30 at\% Ta), where SnO$_{2}$ crystallizes into a tetragonal (cassiterite) phase. Crystallization of SnO$_{2}$ is clearly reflected in the dielectric constant below 30 at\% Ta by two measures: the increased dispersion in the dielectric factor and the marked increase in dissipation factor. For compositions $>$ 30 at\% Ta, Ta-Sn-O remains amorphous for all compositions investigated ($\sim$ 15 at\% Ta $<$ x $<$ $\sim$85 at\% Ta). Nb-Ge-O is x-ray amorphous for $\sim$15 at\% Ta $<$ x $<$ $\sim$ 95 at\% Ta, thus not surprisingly, the dielectric constant and the refractive index scales linearly with composition. It is interesting to note that the dielectric constant in the Nb rich side does not become dispersive since pure Nb$_{2}$O$_{5}$ films are generally conducting. For a more detailed analysis of the dielectric behavior of the oxide, it may be worthwhile to investigate the relaxation mechanism of the oxide as a function of frequency. The density of Ta-Sn-O and Nb-Ge-O varies linearly with composition with the exception of the jump in density for Ta-Sn-O, where SnO$_{2}$ crystallizes.                

There are several different paths we can go down from here. An obvious choice would be to probe the local environment around Ge in Ta-Ge-O using EXAFS. Another low hanging fruit would be to use high energy x-ray diffraction and absorption methods to study Ta-Si-O. Stuctural measurements can be complimented with infrared spectroscopy to directly observe the changes in phonon modes as a function of Ta cation composition for Ta-Ge-O and Ta-Si-O, to establish a more rigorous structure-property relationship. Obviously we are not limited to Ta based systems; we can apply our composition spread approach to rapidly analyze the local structure variation of several different amorphous systems, including ternary oxides. Finally, we can utilize these materials in novel device applications to study changes in device performance with composition. For example, it would be interesting to study the effects of composition in resisitive memory and logic devices such as memristors.
